\subsection{\texorpdfstring{Disintegration of a measure \(\mu\) relative to a \(\mu\) -proper mapping}{Disintegration of a measure mu relative to a mu -proper mapping}}

Let T be a locally compact space having a countable base (in other words, a locally compact Polish space (GT, IX, §6, No.~1). We know that for every positive measure on T, the concepts of integral and essential integral coincide (Ch. V, §1, No.~3, Cor. of Prop. 9). On the other hand, we have the following properties:

Lemma 1. --- If Y is a locally compact space with a countable base, the space \(\mathcal{K}(\mathrm{Y})\) contains a countable dense subset. More precisely, there exists in \(\mathcal{K}(\mathrm{Y})\) a countable subset S consisting of functions \(\geqslant 0\) , such that, for every function \(f \geqslant 0\) of \(\mathcal{K}(\mathrm{Y})\) , there exists a sequence of functions \(f_n \in \mathrm{S}\) ( \(n \geqslant 0\) ) that converges uniformly to \(f\) and is such \(f_n \leqslant f_0\) for all \(n \geqslant 0\) .

For, Y is the union of an increasing sequence \((\mathrm{U}_{n})\) of relatively compact open sets such that \(\overline{U}_{n} \subset U_{n+1}\) for all n (GT, I, §9, No.~9, Prop. 15); the space \(\mathcal{K}(\mathrm{Y})\) is the union of the increasing sequence of Banach spaces \(\mathcal{K}(\mathrm{Y}, \overline{\mathrm{U}}_{n})\) , and each of the latter is known to be separable (GT, Ch. X, §3, No.~3, Th. 1). Let \(S_{n}^{\prime}\) be a countable dense set in \(\mathcal{K}(\mathrm{Y}, \overline{\mathrm{U}}_{n})\) , \(S_{n}\) the set of functions \(\varphi^{+}\) for \(\varphi \in S_{n}^{\prime}\) , and \(u_{n}\) a function in \(\mathcal{K}(\mathrm{Y}, \overline{\mathrm{U}}_{n+1})\) , with values in {[}0,1{]} and equal to 1 on \(U_{n}\) . We take for S the union of the \(S_{n}\) and the set of \(mu_{n}\) for m and n integers \(\geqslant 0\) . For every function \(f \geqslant 0\) of \(\mathcal{K}(\mathrm{Y})\) , f has support contained in one of the \(U_{n}\) , hence is the uniform limit of a sequence of functions \(f_{p} \in S_{n}\) ( \(p \geqslant 1\) ). These functions \(f_{p}\) are uniformly bounded by a positive integer m, and it suffices to take \(f_{0} = mu_{n}\) .

Lemma 2. --- If T is a locally compact space with a countable base, then the Banach space \(\mathcal{K}(\mathrm{Y})\) of continuous numerical functions tending to 0 at the point at infinity is separable.

This lemma is none other than the Cor. to Th. 1 of GT, X, §3, No.~3. One may observe that it also follows from Lemma 1 and the fact that the topology of uniform convergence on \(\mathcal{H}(\mathrm{Y})\) is coarser than the direct limit topology of the topologies of the subspaces \(\mathcal{K}(\mathrm{Y},\overline{\mathrm{U}}_n)\) .

Lemma 3. --- Let T and X be two locally compact spaces with countable bases, \(\mu\) a positive measure on T, and \(t \mapsto \lambda_{t}\) ( \(t \in T\) ) a family of positive measures on X. If the mapping \(t \mapsto \lambda_{t}\) is scalarly \(\mu\) -integrable (for the topology \(\sigma(\mathcal{M}(X), \mathcal{K}(X))\) ), then the family \(t \mapsto \lambda_{t}\) is \(\mu\) -adequate ( \(\S1\) , No.~1, Example).

For, Lemma 1, applied to \(\mathcal{K}(\mathrm{X})\) , shows that the mapping \(t \mapsto \lambda_t\) is vaguely \(\mu\) -measurable (§1, No.~5, Prop. 13).

THEOREM 1. --- Let T and B be two locally compact spaces having countable bases, \(\mu\) a positive measure on T, \(p\) a \(\mu\) -proper mapping (Ch. V, §6, No.~1, Def. 1) of T into B, and \(\nu = p(\mu)\) the image of \(\mu\) under \(p\) . Then there exists a \(\nu\) -adequate family (§1, No.~1, Example) \(b \mapsto \lambda_b\) ( \(b \in B\) ) of positive measures on T, having the following properties:

\begin{enumerate}
\def\labelenumi{\alph{enumi})}
\item
  \(\| \lambda_b\| = 1\) for all \(b\in p(\mathrm{T})\)
\item
  \(\lambda_{b}\) is concentrated on the set \(\overline{p}^{-1}(b)\) (Ch. V, §5, No.~7, Def. 4) for all \(b\in \mathrm{B}\) ; in particular, \(\lambda_{b} = 0\) for \(b\notin p(\mathrm{T})\) ;
\item
  \(\mu = \int \lambda_b d\nu(b)\) .
\end{enumerate}

Moreover, if \(b \mapsto \lambda_b'\) ( \(b \in \mathbf{B}\) ) is a second \(\nu\) -adequate family of positive measures on \(\mathbf{T}\) having the properties \(\mathbf{b}\) ) and \(\mathbf{c}\) ), then \(\lambda_b' = \lambda_b\) almost everywhere in \(\mathbf{B}\) for the measure \(\nu\) .

\begin{enumerate}
\def\labelenumi{\arabic{enumi})}
\tightlist
\item
  Uniqueness. For every function \(f \in \mathcal{K}(\mathrm{B})\) , \(f \circ p\) is \(\mu\) -integrable since \(p\) is \(\mu\) -proper (Ch. V, §6, No.~2, Th. 1); for every function \(g \in \mathcal{K}(\mathrm{T})\) , the function \(t \mapsto g(t)f(p(t))\) is therefore \(\mu\) -integrable. It follows (Ch. V, §3, No.~3, Th. 1) that for almost every \(b \in \mathrm{B}\) , the function \(t \mapsto g(t)f(p(t))\) is \(\lambda_b\) -integrable and that
\end{enumerate}

\[
\int g (t) f \bigl (p (t) \bigr) d \mu (t) = \int d \nu (b) \int g (t) f \bigl (p (t) \bigr) d \lambda_ {b} (t).\tag{1}
\]

But since \(\lambda_{b}\) is concentrated on \(\overline{p}^{-1}(b)\) , we have, for every \(b \in \mathbf{B}\) , \(f(p(t)) = f(b)\) almost everywhere for \(\lambda_{b}\) , therefore the second member of (1) is equal to \(\int f(b)\langle g, \lambda_{b} \rangle d\nu(b)\) . The analogous formula for \(\lambda_{b}'\) also holds; consequently \(\int f(b)\langle g, \lambda_{b} \rangle d\nu(b) = \int f(b)\langle g, \lambda_{b}' \rangle d\nu(b)\) for all \(f \in \mathcal{K}(\mathbf{B})\) and \(g \in \mathcal{K}(\mathbf{T})\) . In other words, the two mappings \(b \mapsto \lambda_{b}\) and \(b \mapsto \lambda_{b}'\) of B into \(\mathcal{M}(T)\) are equal scalarly locally almost everywhere for \(\nu\) , hence equal almost everywhere for \(\nu\) (Lemma 1 and §1, No.~1, Remark 2).

\begin{enumerate}
\def\labelenumi{\arabic{enumi})}
\setcounter{enumi}{1}
\tightlist
\item
  Provisional definition of the family \(b \mapsto \lambda_b\) . For every function \(f \in \mathcal{L}^1(\nu)\) , \(f \circ p\) is \(\mu\) -integrable (Ch. V, §6, No.~2, Th. 1), therefore \((f \circ p) \cdot \mu\) is a bounded measure on \(T\) , and
\end{enumerate}

\[
\left\| (f \circ p) \cdot \mu \right\| = \int | f \circ p | d \mu = \int | f | d \nu = \mathrm{N} _ {1} (f)
\]

(Ch. IV, §4, No.~7, Prop. 12; Ch. V, §5, No.~3, Th. 1 and §6, No.~2, Th. 1). It follows that \((f \circ p) \cdot \mu\) depends only on the class \(\widetilde{f}\) of \(f\) in \(\mathrm{L}^1(\nu)\) and that \(\widetilde{f} \mapsto (f \circ p) \cdot \mu\) is an isometric linear mapping of \(\mathrm{L}^1(\nu)\) into the Banach space \(\mathcal{M}^1(\mathrm{T})\) of bounded measures on \(\mathrm{T}\) , the strong dual of the Banach space \(\overline{\mathcal{K}(\mathrm{T})}\) , which is separable (Lemma 2). By the Dunford-Pettis theorem (§2, No.~5, Cor. 2 of Th. 1) there exists a mapping \(b \mapsto \lambda_b\) of \(\mathrm{B}\) into the unit ball of \(\mathcal{M}^1(\mathrm{T})\) , scalarly \(\nu\) -measurable (for the topology \(\sigma(\mathcal{M}^1(\mathrm{T}), \overline{\mathcal{K}(\mathrm{T})})\) ) and such that, for every function \(f \in \mathcal{L}^1(\nu)\) ,

\[
(f \circ p) \cdot \mu = \int f (b) \lambda_ {b} d \nu (b),\tag{2}
\]

which may also be written, for every function \(g \in \overline{\mathcal{K}(\mathrm{T})}\)

\[
\int g (t) f (p (t)) d \mu (t) = \int f (b) d \nu (b) \int g (t) d \lambda_ {b} (t).\tag{3}
\]

If \(f \geqslant 0\) and \(g \geqslant 0\) , the first member of (3) is \(\geqslant 0\) , which proves that for every function \(g \geqslant 0\) in \(\mathcal{K}(\mathrm{T})\) , the measure \(\left( \int g(t) d\lambda_b(t) \right) \cdot \nu\) is \(\geqslant 0\) , hence that \(\int g(t) d\lambda_b(t) \geqslant 0\) except for \(b\) belonging to a \(\nu\) -negligible set \(\mathrm{N}(g)\) (Ch. V, §5, No.~3, Cor. 3 of Prop. 3). Now, there exists a dense sequence \((g_n)\) in the space \(\mathcal{K}_+(\mathrm{T})\) of functions \(\geqslant 0\) of \(\mathcal{K}(\mathrm{T})\) (Lemma 1). The union \(\mathrm{N}\) of the \(\mathrm{N}(g_n)\) is \(\nu\) -negligible and, for \(b \notin \mathrm{N}\) , we have \(\int g_n(t) d\lambda_b(t) \geqslant 0\) for all \(n\) , therefore \(\int g(t) d\lambda_b(t) \geqslant 0\) for every function \(g \in \mathcal{K}_+(\mathrm{T})\) , in other words \(\lambda_b \geqslant 0\) .

This being so, we may replace \(\lambda_{b}\) by 0 for every \(b \in N\) without altering the validity of (3); we can therefore assume this modification to have been made, so that \(\lambda_{b} \geqslant 0\) for every \(b \in B\) .

\begin{enumerate}
\def\labelenumi{\arabic{enumi})}
\setcounter{enumi}{2}
\tightlist
\item
  Extensions of the formula (3).
\end{enumerate}

\(\alpha)\) For every function \(f \in \mathcal{L}^1(\nu)\) , it follows from (3) that the mapping \(b \mapsto \lambda_b\) of B into \(\mathcal{M}(T)\) is scalarly integrable for the measure \(|f \cdot \nu|\) and the topology \(\sigma(\mathcal{M}(T), \mathcal{K}(T))\) , therefore (Lemma 3) the family \(b \mapsto \lambda_b\) is \(|f \cdot \nu|\) -adequate. Now let \(g\) be a numerical function defined on \(T\) , integrable for the measure \(|(f \circ p) \cdot \mu|\) , that is (Ch. V, §5, No.~3, Th. 1), such that \(t \mapsto g(t)f((p(t))\) is \(\mu\) -integrable; it then follows from (2), from Th. 1 of Ch. V, §3, No.~3 and from Th. 1 of Ch. V, §5, No.~3 that, for almost every \(b \in \mathbf{B}\) , \(g\) is integrable for \(\lambda_b\) , that the function (defined almost everywhere) \(b \mapsto \int g(t)d\lambda_b(t)\) is integrable for \(|f \cdot \nu|\) , and that the formula (3) is again valid.

\(\beta)\) For every function \(g \in \overline{\mathcal{K}(\mathrm{T})}\) , it follows from (3), applied to \(f \in \mathcal{K}(\mathrm{B})\) , that the mapping \(p\) is proper for the measure \(|g \cdot \mu|\) (Ch. V, §6, No.~1, Def. 1) and the image under \(p\) of the measure \(g \cdot \mu\) is the measure with density \(b \mapsto \int g(t) d\lambda_b(t)\) with respect to \(\nu\) . If \(f\) is then taken to be a function such that \(f \circ p\) is integrable for the measure \(|g \cdot \mu|\) , that is, such that \(t \mapsto g(t)f(p(t))\) is \(\mu\) -integrable (Ch. V, §5, No.~3, Th. 1), the formula (3) is again valid (Ch. V, §6, No.~2, Th. 1).

\begin{enumerate}
\def\labelenumi{\arabic{enumi})}
\setcounter{enumi}{3}
\tightlist
\item
  Properties of the family \(b \mapsto \lambda_b\) . By the property \(\beta\) ), we can apply formula (3) by taking \(f = 1\) , \(g \in \mathcal{K}(\mathrm{T})\) ; this proves that \(b \mapsto \lambda_b\) is scalarly \(\nu\) -integrable (for the topology \(\sigma(\mathcal{M}(\mathrm{T}), \mathcal{K}(\mathrm{T}))\) ), hence is \(\nu\) -adequate (Lemma 3), and that \(\mu = \int \lambda_b d\nu(b)\) .
\end{enumerate}

Now let \(\psi\) be any function in \(\mathcal{K}(\mathrm{B})\) ; the conditions of property \(\alpha\) are fulfilled by taking \(f \in \mathcal{K}(\mathrm{B})\) and \(g = \psi \circ p\) , because the function \(\psi(p(t))f(p(t))\) is \(\mu\) -integrable since \(f\psi\) belongs to \(\mathcal{K}(\mathrm{B})\) and \(p\) is \(\mu\) -proper. Then \(\psi \circ p\) is \(\lambda_b\) -integrable for almost every \(b \in \mathrm{B}\) , and

\[
\int f (p (t)) \psi (p (t)) d \mu (t) = \int f (b) d \nu (b) \int \psi (p (t)) d \lambda_ {b} (t);
\]

but the first member is by definition \(\int f(b)\psi (b)d\nu (b)\) . We therefore see that for every function \(\psi \in \mathcal{K}(\mathrm{B})\) , the measure \(\psi \cdot \nu\) and the measure with density \(b\mapsto \int \psi (p(t))d\lambda_b(t)\) are identical. Consequently (Ch. V, §5, No.~3, Cor. 2 of Prop. 3) there exists a \(\nu\) -negligible set \(\mathrm{N}'(\psi)\) such that, for every \(b\notin \mathrm{N}'(\psi)\) , the function \(\psi \circ p\) is \(\lambda_{b}\) -integrable and \(\psi (b) = \int \psi (p(t))d\lambda_b(t)\) .

Let S be a countable subset of \(\mathcal{K}(\mathrm{B})\) having the properties stated in Lemma 1 (with Y = B), and let N' be the \(\nu\) -negligible set that is the union of the N'(ψ) for \(\psi \in \mathrm{S}\) . Every function \(\psi \geqslant 0\) of \(\mathcal{K}(\mathrm{B})\) is the uniform limit of a sequence \((\psi_n)\) of elements of S with \(\psi_n \leqslant \psi_0\) . Consequently for \(b \notin \mathrm{N}'\) , Lebesgue's theorem shows on the one hand that \(\psi \circ p\) is \(\lambda_b\) -integrable, in other words that \(p\) is \(\lambda_b\) -proper, and on the other hand that \(\psi(b) = \int \psi(p(t)) d\lambda_b(t)\) . In other terms, the mappings \(b \mapsto \varepsilon_b\) and \(b \mapsto p(\lambda_b)\) of B into \(\mathcal{M}(\mathrm{B})\) (the latter being defined almost everywhere) are scalarly almost everywhere equal for \(\nu\) (and for the topology \(\sigma(\mathcal{M}(\mathrm{B}), \mathcal{K}(\mathrm{B}))\) ); it follows that these mappings are equal almost everywhere for \(\nu\) (Lemma 1 and §1, No.~1, Remark 2). Finally, if \(p(\lambda_b) = \varepsilon_b\) , the set B - \{b\} is \(\varepsilon_b\) -negligible, therefore the set T - \(\overline{p}^{-1}(b)\) is \(\lambda_b\) -negligible (Ch. V, §6, No.~2, Cor. 2 of Prop. 2), in other words \(\lambda_b\) is concentrated on \(\overline{p}^{-1} (b)\) ; and, on the other hand, \(\| \lambda_b\| = \int d\lambda_b = \int d((p(\lambda_b)) = \| \varepsilon_b\| = 1\) (Ch. V, §6, No.~2, Th. 1).

\begin{enumerate}
\def\labelenumi{\arabic{enumi})}
\setcounter{enumi}{4}
\tightlist
\item
  Modifications of the family \(b \mapsto \lambda_{b}\) . We have thus defined a \(\nu\) -adequate family \(b \mapsto \lambda_{b}\) of measures \(\geqslant 0\) on T, satisfying condition c) of the statement and such that, for almost every \(b \in B\) , p is \(\lambda_{b}\) -proper, and \(\lambda_{b}\) is concentrated on \(\overline{p}^{-1}(b)\) and has norm 1. Let \(N^{\prime\prime}\) be the \(\nu\) -negligible set of points \(b \in B\) where one of the last three properties is not verified; we can then modify \(\lambda_{b}\) in the following manner. If \(b \in B - p(T)\) , take \(\lambda_{b} = 0\) ; if \(b \in p(T) \cap N^{\prime\prime}\) , take \(\lambda_{b} = \varepsilon_{\xi(b)}\) , where \(\xi(b)\) is any point of \(\overline{p}^{-1}(b)\) . Since \(B - p(T)\) is \(\nu\) -negligible (Ch. V, §6, No.~2, Cor. 3 of Prop. 2), we have only modified \(\lambda_{b}\) at the points of a negligible set, consequently the family \(b \mapsto \lambda_{b}\) is still \(\nu\) -adequate and has the property c); moreover, it now satisfies a) and b), which completes the proof.
\end{enumerate}

Every \(\nu\) -adequate family \(b \mapsto \lambda_{b}\) of positive measures on T, having the properties b) and c) of Th. 1, is said to be a disintegration of the measure \(\mu\) , relative to the \(\mu\) -proper mapping p.

